\documentclass[10pt,a4paper]{amsart}
\usepackage[margin=19mm]{geometry}
\usepackage{amsmath,amssymb,mathtools}
\usepackage[scr=rsfs,cal=euler]{mathalfa}
\usepackage{xfrac}
\usepackage[unicode,hidelinks,bookmarks=false]{hyperref}
\newcommand{\ie}{i.e.\ }
\newcommand{\cf}{cf.\ }
\newcommand{\resp}{respectively\ }
\newcommand{\Subsection}{\subsection}
\providecommand{\nobreakdash}{\nobreak\hbox{-}}
\setlength{\emergencystretch}{2em}
\multlinegap0pt
\usepackage{bm}
\usepackage{bbm}
\usepackage{dsfont}
\usepackage{xspace}
\usepackage{cancel}
\usepackage{mathtools}
\usepackage{amsthm}
\makeatletter
\@ifundefined{theorem}{\newtheorem{theorem}{Theorem}}{}
\@ifundefined{corollary}{\newtheorem{corollary}[theorem]{Corollary}}{}
\@ifundefined{lemma}{\newtheorem{lemma}[theorem]{Lemma}}{}
\makeatother
\title[An exact Ramsey number of large bipartite graphs versus odd ]{An exact Ramsey number of large bipartite graphs versus odd wheel}
\author{Sayan Gupta, Kaushik Majumder}
\date{Source: arXiv:2511.14867v2 (CC BY 4.0)}
\begin{document}
\maketitle

\noindent\emph{Source and licence.} An exact Ramsey number of large bipartite graphs versus odd wheel. Version arXiv:2511.14867v2 (CC BY 4.0), distributed under the Creative Commons Attribution 4.0 International licence, as recorded in the arXiv licence metadata for this exact version and shown on the abstract page https://arxiv.org/abs/2511.14867v2. Full rights notice: licenses/arxiv-2511.14867-CC-BY-4.0.txt.\par\smallskip

\noindent\textit{Editorial scope.} Counted unit: the Lemma whose source label is \texttt{connectivity-first range} in the article (source0.tex lines 456--458, source label \texttt{connectivity-first range}), delivered together with the article's own complete original proof of it (source0.tex lines 459--548, one \texttt{proof} environment of 90 lines). No mathematics is written, repaired or re-derived: statement and proof are the article's own text line for line. Required context retained uncounted: nbd dependency (lines 110--112); nbd complementary dependency (lines 167--169); K\_2n freeness dependency (lines 171--173); intersection lemma 2 (lines 277--282); upper bound Del(barG) (lines 290--292). Every definition, display and label of those lines is retained unchanged. Omitted from this package and not redistributed: 1--109, 113--166, 170--170, 174--276, 283--289, 293--455, 549--1085. Nothing inside a retained excerpt is omitted or altered: every retained body is delivered exactly as the article's lines are written, so no omission receipt of any kind accompanies this package. Citations: the retained excerpts carry no citation command at all, so no bibliography entry of the article is transcribed and no reference is redistributed. Reference adaptation: none was needed and none was made; every reference inside the delivered unit resolves inside it, so no unresolved reference or citation is produced. Printed slips kept literal: no printed word, symbol, hypothesis or number of the counted statement or of its proof was altered. Layout and class: the article's own class is \texttt{amsart} with packages including amsfonts, amssymb, amsthm, amsbsy, latexsym, amscd, amsmath, dsfont, euscript, enumerate, geometry, graphicx, ifthen, manfnt; the wrapper is the lane's pinned \texttt{amsart} editorial wrapper at 10pt on A4, which restates the theorem-like environments the retained text needs and adds the article's own macro definitions that the retained text needs (none), with extra wrapper packages (bm, bbm, dsfont, xspace, cancel, mathtools). The delivered numbering is the unit's own. Assets: no retained span contains a figure environment, a graphics-inclusion command, a PDF-page inclusion, a table or a TikZ picture; no image file is copied into this package, the package declares no asset dependency and no image file is redistributed with it. Licence: version arXiv:2511.14867v2 of this article is distributed under the Creative Commons Attribution 4.0 International licence, as recorded in the arXiv licence metadata for this exact version and shown on the abstract page https://arxiv.org/abs/2511.14867v2. A full rights notice is delivered at licenses/arxiv-2511.14867-CC-BY-4.0.txt. Characters: every retained excerpt is pure ASCII, so the delivered engine, XeLaTeX with its default Latin Modern text font, reports no missing character.
\begin{equation}\label{nbd dependency}
N_{G[Q]}(x)\subset N_{G}(x)\subset N_{G[Q]}(x)\sqcup P.
\end{equation}

\begin{equation}\label{nbd complementary dependency}
\{u\}\sqcup N_{G}(u)\sqcup N_{\overline{G}}(u)=V(G),
\end{equation}

\begin{equation}\label{K_2n freeness dependency}
|N_{G}(u)\cap N_{G}(v)|\leq n-1
\end{equation}

\begin{corollary}\label{intersection lemma 2}
Let $G$ be a graph with $A\subset B=V(G)$ and $|N_{G}(a)|\geq d$ for each $a\in A$. Then there exist at least two distinct vertices $x,y\in A$ such that 
\begin{equation*}
 |N_{G}(x)\cap N_{G}(y)|> \frac{d^{2}}{|B|}-\frac{d}{|A|}=\frac{d^{2}}{|V(G)|}-\frac{d}{|A|}.   
\end{equation*}
\end{corollary}

\begin{lemma}\label{upper bound Del(barG)}
 $\Delta(\overline{G})\leq 2n+2=2(n+1)$. 
\end{lemma}

\begin{lemma}\label{connectivity-first range}
Within this range, along with the vertex set decomposition $P\sqcup Q=V(\overline{H})$. Then the connectivity parameter of $\overline{G}[Q]$ satisfies $\kappa(\overline{G}[Q])\geq2$.
\end{lemma}

\begin{proof}
If not, we assume $\kappa(\overline{G}[Q])\in\{0,1\}$ and 
construct	
\begin{equation*}
	R=\begin{cases}
		\emptyset &\text{ if } \kappa(\overline{G}[Q])=0\\
		\{q\}&\text{ if } \kappa(\overline{G}[Q])=1\text{ i.e. else},
	\end{cases}
\end{equation*}	
where, in the second case, $q$ is a \emph{cut-vertex} of $\overline{G}[Q]$. By construction, the induced subgraph  $\overline{G}[Q-R]$ is not connected.  

On $Q-R$, we construct the relation $\sim$ by declaring $x\sim y$ if and only if there is a path in $\overline{G}[Q-R]$ from 
$x$ to $y$. (We allow the  $2-$tuples $(x,x)$ as paths, where $x\in Q-R$, for the reflexive property.) This is an equivalence relation, and its equivalent classes are precisely the maximal connected components of $\overline{G}[Q-R]$. We call these classes $S_{1},\ldots,S_{m}$, where $m\in[2,9]\cap\mathbb{Z}$, so that 
\begin{equation*}
Q=R\sqcup S_{1}\sqcup\ldots\sqcup S_{m}.
\end{equation*}
Due to the maximal property, for each  $S,S'\in\{S_{1},\ldots,S_{m}\}$ with $S\neq S'$, there is no edge of the form $\{s,s'\}$, where $s\in S$ and $s'\in S'$. Due to these properties, the neighbourhood relation \eqref{nbd dependency} gets a revision, that is, for each $x\in S$ , where $S\in\{S_{1},\ldots,S_{m}\}$
\begin{equation*}
 N_{\overline{G}[S]}(x)\subset N_{\overline{G}[Q]}(x)\subset N_{\overline{G}[S]}(x)\sqcup R.    
\end{equation*}
Since each $S\in\{S_{1},\ldots,S_{m}\}$ forms a connected component, thus we have 
\begin{equation*}
 \frac{V(\overline{H})}{10}+1\leq\delta(\overline{G}[Q])\leq\delta(\overline{G}[S])\leq|S|-1.   
\end{equation*}
This means $|S|\geq2+\frac{V(\overline{H})}{10}$ for each $S\in\{S_{1},\ldots,S_{m}\}$. Since 
\begin{equation*}
|V(\overline{H})|-|P|-|R|=|Q|-|R|=|S_{1}|+\ldots+|S_{m}|\geq2m+m\frac{V(\overline{H})}{10}    
\end{equation*}
holds we have $m\leq9$. 

We choose $B\in\{S_{1},\ldots,S_{m}\}$ such that $|B|=\min\{|S_{1}|,\ldots,|S_{m}|\}$ and construct the set 
\begin{equation*}
 A=(S_{1}\sqcup\ldots\sqcup S_{m})-B.   
\end{equation*}

\noindent{\textsl{Claim} :}  $|A|\in\left[\frac{|V(\overline{H})|}{2}-1,n-1\right]\cap\mathbb{Z}$.
\begin{proof}[\tt{Proof of the claim} :]\renewcommand{\qedsymbol}{}
Here $|A|\geq|B|$. Consequently, $|A|\geq\frac{|Q|-|R|}{2}=\frac{|V(\overline{H})|-|P|-|R|}{2}\geq|B|$. Thus using $|P|\leq1$ and $|R|\leq1$, we have $|A|\geq\frac{|V(\overline{H})|}{2}-1$. We also have $|A|\leq n-1$. If not, let $|A|\geq n$. In the complementary graph of $\overline{G}$, i.e. in $G$, for each $a\in A$ and $b\in B$, $\{a,b\}$ forms an edge in $G$. We choose two vertices $b,b'\in B$, then $N_{G}(b)\cap A=A=N_{G}(b')\cap A$. This implies 
\begin{equation*}
 |N_{G}(b)\cap N_{G}(b')|\geq|N_{G}(b)\cap N_{G}(b')\cap A|=|A|\geq n,   
\end{equation*}
which contradicts \eqref{K_2n freeness dependency} and establishes the claim. 
\end{proof}

We  consider the vertex partition $A\sqcup C\sqcup\{v\}=V(G)$ and construct the induced subgraphs 
\begin{equation*}
 G[A\sqcup C],\quad\overline{G}[A\sqcup C]\quad\text{ and }\quad G[C].   
\end{equation*}
Here $B\subset C$.  For a vertex $w\in B$, since $v$ is the maximum degree vertex in $\overline{G}$ and $w$ is adjacent to $v$ in $\overline{G}$, using \eqref{nbd dependency}, we have 
\begin{equation*}
 N_{\overline{G}}(w)=N_{\overline{G}[A\sqcup C]}(w)\sqcup\{v\}.   
\end{equation*}
This implies, for each $w\in B$,
\begin{equation*}
|N_{\overline{G}[A\sqcup C]}(w)|\leq\Delta(\overline{G})-1=|V(\overline{H})|-1.
\end{equation*}
Using \eqref{nbd complementary dependency} and Lemma~\ref{upper bound Del(barG)}, we have the following:- 
\begin{align*}
|N_{G[A\sqcup C]}(w)|&=(|A|+|C|)-1-|N_{\overline{G}[A\sqcup C]}(w)|\\
&=(|V(G)|-1)-1-|N_{\overline{G}[A\sqcup C]}(w)|\\
&\geq |V(G)|-|V(\overline{H})|-1.
\end{align*}
Revising \eqref{nbd dependency}, we have $N_{G[A\sqcup C]}(b)=N_{G[A]}(b)\sqcup N_{G[C]}(b)$ for each $b\in B$. Therefore for each $b\in B$,
\begin{align*}
|N_{G[C]}(b)|=|N_{G[A\sqcup C]}(b)|-|N_{G[A]}(b)|\geq &|N_{G[A\sqcup C]}(b)|-|A|\\
&\geq |V(G)|-|V(\overline{H})|-|A|-1=|C|-|V(\overline{H})|.   
\end{align*}
Using Corollary~\ref{intersection lemma 2} over the graph $G[C]$, with $d=|C|-|V(\overline{H})|$ and $B\subset C$,  we conclude that there exist two distinct $b,b'\in B$ such that
\begin{equation*}
|N_{G[C]}(b)\cap N_{G[C]}(b')|>\frac{d^{2}}{|C|}-\frac{d}{|B|}\geq n-1-|A|+\frac{P(n,\varepsilon)}{Q(n,\varepsilon)}\geq n-1-|A|,    
\end{equation*}
where 
\begin{align*}
P(n,\varepsilon)&=(15+12\varepsilon)(n+1)^{2}+(16\varepsilon^{2}-36\varepsilon+32)(n+1)-4(4\varepsilon^{3}+\varepsilon^{2}-4)\\
Q(n,\varepsilon)&=[9(n+1)+(4+2\varepsilon)][3(n+1)+(4-2\varepsilon)]. 
\end{align*}
For each $n\geq 1$ and $\varepsilon\in\left(1,\frac{(2\sqrt{3}-3)(n+1)}{2}\right]$, $\frac{P(n,\varepsilon)}{Q(n,\varepsilon)}\geq0$, concludes the last inequality. While estimating the quantity $\frac{d^{2}}{|C|}-\frac{d}{|B|}$, we have used the following inequalities:- 
\begin{align*}
\frac{|V(\overline{H})|^{2}}{|C|}\geq&\frac{|V(\overline{H})|^{2}}{|V(G)|-\frac{|V(\overline{H})|}{2}}=\frac{[3(n+1)-2\varepsilon]^{2}}{9(n+1)+(4+2\varepsilon)},\\
\frac{|V(\overline{H})|-|C|}{|B|}\geq&\frac{\frac{3}{2}|V(\overline{H})|-|V(G)|}{\frac{|V(\overline{H})|}{2}+1}=3-\frac{|V(G)|+3}{\frac{|V(\overline{H})|}{2}+1}=3-\frac{12(n+1)+16}{3(n+1)+(4-2\varepsilon)}.
\end{align*}
Here $N_{G}(b)\cap N_{G}(b')\cap A=A$, $v\notin N_{G}(b)\cap N_{G}(b')$, $C\subset V(G)$ and $G[C]$ forms a subgraph of $G$, these imply
for two distinct vertices $b,b'\in B$,
\begin{align*}
|N_{G}(b)\cap N_{G}(b')|=&|N_{G}(b)\cap N_{G}(b')\cap A|+|N_{G}(b)\cap N_{G}(b')\cap C|\\
=&|N_{G}(b)\cap N_{G}(b')\cap A|+|N_{G[C]}(b)\cap N_{G[C]}(b')|\\
&\geq |A|+(n-1)-|A|+1=n.
\end{align*}
This contradicts \eqref{K_2n freeness dependency}.
\end{proof}

\end{document}
